\chapter{Описание стенда и подготовка}
\label{ch:stand}

\section{Конфигурация виртуального стенда}

Виртуальная машина \texttt{gateway} --- Ubuntu~22.04~LTS (VirtualBox),
два сетевых адаптера:
\begin{itemize}
  \item \textbf{Адаптер~1} (\texttt{enp0s3}) --- тип \textit{NAT}:
        выход в интернет, IP \texttt{10.0.2.15/24}, шлюз \texttt{10.0.2.2};
  \item \textbf{Адаптер~2} (\texttt{enp0s8}) --- тип \textit{Internal Network
        intnet}: LAN-сегмент, IP \texttt{192.168.\cfgLabStudentN.1/24}.
\end{itemize}

Конфигурация выполнена в лаб.~работах~4 и~5: маршрутизация, DHCP, DNS (BIND9).

\vspace{6pt}
Виртуальная машина \texttt{mail} --- Ubuntu~22.04~LTS:
\begin{itemize}
  \item \textbf{Адаптер~1} (\texttt{enp0s3}) --- тип \textit{Internal Network
        intnet}, IP \texttt{192.168.\cfgLabStudentN.5/24},
        шлюз \texttt{192.168.\cfgLabStudentN.1},
        DNS \texttt{192.168.\cfgLabStudentN.1}.
\end{itemize}

Объём RAM --- не менее 2~ГБ; объём диска --- не менее 20~ГБ (iRedMail
требует место под почтовые ящики \texttt{/var/vmail}).

% --- Скриншот: настройки ВМ mail в VirtualBox (вкладка Сеть) ---
\begin{figure}[h]
  \centering
  \includegraphics[width=0.85\textwidth]{img/01_vbox_mail_settings}
  \caption{Настройки ВМ \texttt{mail} --- адаптер в VirtualBox}
  \label{fig:vbox_mail}
\end{figure}

\section{Настройка статического IP (mail)}

На ВМ \texttt{mail} редактируем файл netplan:
\begin{lstlisting}
nano /etc/netplan/01-netcfg.yaml
\end{lstlisting}

Содержимое:
\begin{lstlisting}
network:
  version: 2
  renderer: networkd
  ethernets:
    enp0s3:
      dhcp4: no
      addresses: [192.168.N.5/24]
      gateway4: 192.168.N.1
      nameservers:
        addresses: [192.168.N.1]
        search: [yazikov.iks531.local]
\end{lstlisting}

Применяем конфигурацию:
\begin{lstlisting}
netplan apply
\end{lstlisting}

% --- Скриншот: nano /etc/netplan/01-netcfg.yaml ---
\begin{figure}[h]
  \centering
  \includegraphics[width=0.85\textwidth]{img/02_netplan_mail}
  \caption{Файл \texttt{netplan} ВМ \texttt{mail} --- статический IP}
  \label{fig:netplan_mail}
\end{figure}

\section{Назначение FQDN хосту mail}

Для корректной работы iRedMail хост должен иметь полное доменное имя.
Выполняем команды:
\begin{lstlisting}
hostnamectl set-hostname mail.yazikov.iks531.local
\end{lstlisting}

Обновляем \texttt{/etc/hosts}:
\begin{lstlisting}
nano /etc/hosts
\end{lstlisting}

Содержимое:
\begin{lstlisting}
127.0.0.1  mail.yazikov.iks531.local mail localhost
192.168.N.5  mail.yazikov.iks531.local mail
\end{lstlisting}

Проверка:
\begin{lstlisting}
hostname -f
\end{lstlisting}
Должен вернуть \texttt{mail.yazikov.iks531.local}.

% --- Скриншот: hostname -f и cat /etc/hosts ---
\begin{figure}[h]
  \centering
  \includegraphics[width=0.85\textwidth]{img/03_hostname_fqdn}
  \caption{FQDN ВМ \texttt{mail} --- вывод \texttt{hostname -f}}
  \label{fig:hostname_fqdn}
\end{figure}

\section{Добавление DNS-записи для mail на gateway}

Переходим на ВМ \texttt{gateway}. Добавляем A-запись и PTR-запись
запуском подготовленного скрипта:
\begin{lstlisting}
sudo bash gateway_add_mail_dns.sh
\end{lstlisting}

Либо вручную --- редактируем прямую зону:
\begin{lstlisting}
nano /var/lib/bind/forward.db
\end{lstlisting}

Добавляем строку:
\begin{lstlisting}
mail    IN      A       192.168.N.5
\end{lstlisting}

Перезапускаем BIND9:
\begin{lstlisting}
systemctl restart bind9
\end{lstlisting}

Проверяем разрешение имени с ВМ \texttt{mail}:
\begin{lstlisting}
nslookup mail
nslookup 192.168.N.5
\end{lstlisting}

% --- Скриншот: вывод nslookup mail (успешное разрешение) ---
\begin{figure}[h]
  \centering
  \includegraphics[width=0.85\textwidth]{img/04_nslookup_mail}
  \caption{DNS-разрешение имени \texttt{mail} с ВМ \texttt{mail}}
  \label{fig:nslookup_mail}
\end{figure}

\section{Проверка доступа в интернет}

Перед установкой iRedMail необходимо убедиться, что ВМ \texttt{mail}
имеет выход в интернет для скачивания пакетов:
\begin{lstlisting}
ping -c3 ya.ru
apt-get update
\end{lstlisting}

% --- Скриншот: ping ya.ru с ВМ mail ---
\begin{figure}[h]
  \centering
  \includegraphics[width=0.85\textwidth]{img/05_ping_internet}
  \caption{Проверка доступа в интернет с ВМ \texttt{mail}}
  \label{fig:ping_inet}
\end{figure}

\endinput
